\documentclass[10pt,a4paper]{amsart}
\usepackage[margin=19mm]{geometry}
\usepackage{amsmath,amssymb,mathtools}
\usepackage[scr=rsfs,cal=euler]{mathalfa}
\usepackage{xfrac}
\usepackage[unicode,hidelinks,bookmarks=false]{hyperref}
\newcommand{\ie}{i.e.\ }
\newcommand{\cf}{cf.\ }
\newcommand{\resp}{respectively\ }
\newcommand{\Subsection}{\subsection}
\providecommand{\nobreakdash}{\nobreak\hbox{-}}
\setlength{\emergencystretch}{2em}
\multlinegap0pt
\usepackage{bm}
\usepackage{bbm}
\usepackage{dsfont}
\usepackage{xspace}
\usepackage{cancel}
\usepackage{mathtools}
\usepackage{amsthm}
\makeatletter
\@ifundefined{proposition}{\newtheorem{proposition}{Proposition}}{}
\makeatother
\newcommand{\RR}{\mathbb{R}}
\title[On the geometry of weak convergence without total variation ]{On the geometry of weak convergence without total variation convergence}
\author{Nicola Bariletto, Stephen G. Walker}
\date{Source: arXiv:2608.03290v1 (CC BY 4.0)}
\begin{document}
\maketitle

\noindent\emph{Source and licence.} On the geometry of weak convergence without total variation convergence. Version arXiv:2608.03290v1 (CC BY 4.0), distributed under the Creative Commons Attribution 4.0 International licence, as recorded in the arXiv licence metadata for this exact version and shown on the abstract page https://arxiv.org/abs/2608.03290v1. Full rights notice: licenses/arxiv-2608.03290-CC-BY-4.0.txt.\par\smallskip

\noindent\textit{Editorial scope.} Counted unit: the Proposition whose source label is \texttt{prop:weak} in the article (source0.tex lines 738--740, source label \texttt{prop:weak}), delivered together with the article's own complete original proof of it (source0.tex lines 741--798, one \texttt{proof} environment of 58 lines). No mathematics is written, repaired or re-derived: statement and proof are the article's own text line for line. Required context retained uncounted: none. Every definition, display and label of those lines is retained unchanged. Omitted from this package and not redistributed: 1--737, 799--908. Nothing inside a retained excerpt is omitted or altered: every retained body is delivered exactly as the article's lines are written, so no omission receipt of any kind accompanies this package. Citations: the retained excerpts carry no citation command at all, so no bibliography entry of the article is transcribed and no reference is redistributed. Reference adaptation: none was needed and none was made; every reference inside the delivered unit resolves inside it, so no unresolved reference or citation is produced. Printed slips kept literal: no printed word, symbol, hypothesis or number of the counted statement or of its proof was altered. Layout and class: the article's own class is \texttt{article} with packages including geometry, amsmath, amssymb, amsthm, amsfonts, mathrsfs, graphicx, xcolor, natbib, authblk; the wrapper is the lane's pinned \texttt{amsart} editorial wrapper at 10pt on A4, which restates the theorem-like environments the retained text needs and adds the article's own macro definitions that the retained text needs (\texttt{\textbackslash RR}), with extra wrapper packages (bm, bbm, dsfont, xspace, cancel, mathtools). The delivered numbering is the unit's own. Assets: no retained span contains a figure environment, a graphics-inclusion command, a PDF-page inclusion, a table or a TikZ picture; no image file is copied into this package, the package declares no asset dependency and no image file is redistributed with it. Licence: version arXiv:2608.03290v1 of this article is distributed under the Creative Commons Attribution 4.0 International licence, as recorded in the arXiv licence metadata for this exact version and shown on the abstract page https://arxiv.org/abs/2608.03290v1. A full rights notice is delivered at licenses/arxiv-2608.03290-CC-BY-4.0.txt. Characters: every retained excerpt is pure ASCII, so the delivered engine, XeLaTeX with its default Latin Modern text font, reports no missing character.
\begin{proposition}\label{prop:weak}
The sequence of probability measures $F_j$ converges to $G$ in L\'evy--Prokhorov distance.
\end{proposition}

\begin{proof}
Since $d_{w}$ metrizes weak convergence, it suffices to show that for any bounded, continuous function $\phi \in C_b(\RR^2)$,
\begin{equation*}
    \lim_{j \to \infty} \int_{\RR^2} f_j(x,y) \phi(x,y) \, dx \, dy = \int_{\RR^2} g(x,y)\phi(x,y) \, dx \, dy.
\end{equation*}
Because both $f_j$ and $g$ evaluate to zero outside of $\Omega$, the integrals may be restricted
to $\Omega$ without loss of generality. First, we analyze the weak limit of the indicator function $1_{A_j}$. Write
\begin{equation*}
    \int_\Omega 1_{A_j} \phi(x,y) \, dx \, dy = \int_{T_j} \phi(x,y) \, dx \, dy + \int_{B_j \setminus T_j} \phi(x,y) \, dx \, dy.
\end{equation*}
Since $\phi$ is bounded, let $M := \sup_{(x,y)\in\RR^2} |\phi(x,y)|<\infty$, so that
\begin{equation*}
    \left| \int_{B_j \setminus T_j} \phi(x,y) \, dx \, dy \right| \leq M \lambda_2(B_j) = \frac{M}{3j} \to 0
\end{equation*}
as $j\to\infty$. Now define the marginal integral $\Phi(x) := \int_0^1 \phi(x,y) \, dy$, so that Fubini's theorem yields
\begin{equation*}
    \int_{T_j} \phi(x,y) \, dx \, dy = \sum_{k=0}^{j-1} \int_{k/j}^{k/j + 1/(3j)} \Phi(x) \, dx.
\end{equation*}
Since $\phi$ is continuous on the compact set $\Omega$, it is uniformly continuous
there, with some modulus of continuity $\omega_\phi$; that is, $|\phi(z) - \phi(z')| \leq
\omega_\phi(|z - z'|)$ for all $z, z' \in \Omega$, with $\omega_\phi(\delta) \to 0$ as
$\delta \to 0^+$. The marginal $\Phi$ inherits this modulus, since for $x, x' \in [0,1]$
\begin{equation*}
    |\Phi(x) - \Phi(x')|
    = \left| \int_0^1 \bigl( \phi(x,y) - \phi(x',y) \bigr) dy \right|
    \leq \int_0^1 \bigl| \phi(x,y) - \phi(x',y) \bigr| \, dy
    \leq \omega_\phi(|x - x'|),
\end{equation*}
the last step using the fact that $|(x,y) - (x',y)| = |x - x'|$. Thus $\Phi$ is uniformly continuous on $[0,1]$
with modulus $\omega_\phi$. On each sub-interval $[k/j,\, k/j + 1/(3j)]$, every point lies within
$1/(3j)$ of the left endpoint $k/j$, so that
\begin{equation*}
    \left| \int_{k/j}^{k/j + 1/(3j)} \Phi(x)\,dx - \frac{1}{3j}\,\Phi\!\left(\frac{k}{j}\right) \right|
    = \left| \int_{k/j}^{k/j + 1/(3j)} \left( \Phi(x) - \Phi\!\left(\frac{k}{j}\right) \right) dx \right|
    \leq \frac{1}{3j}\,\omega_\phi\!\left(\frac{1}{3j}\right).
\end{equation*}
Summing over the $j$ sub-intervals, the total error is bounded by
$\frac{1}{3}\,\omega_\phi(1/(3j)) = o(1)$, whence
\begin{equation*}
    \int_{T_j} \phi(x,y) \, dx \, dy = \frac{1}{3} \sum_{k=0}^{j-1}\frac{1}{j}  \Phi\left(\frac{k}{j}\right) + o(1).
\end{equation*}
The first term is one-third of a Riemann sum for $\Phi$ over $[0,1]$, so that
\begin{equation*}
    \lim_{j \to \infty} \int_{T_j} \phi(x,y) \, dx \, dy = \frac{1}{3} \int_0^1 \Phi(x) \, dx = \frac{1}{3} \int_\Omega \phi(x,y) \, dx \, dy.
\end{equation*}
Now, substituting $1_{\Omega \setminus A_j} = 1_\Omega - 1_{A_j}$ into the definition of $f_j$
yields
\begin{equation*}
    f_j = \left( \left( 1 - \frac{p_j}{1-p_j} \right) 1_\Omega + \left( 1 + \frac{p_j}{1-p_j} \right) 1_{A_j} \right) 1_\Omega = \frac{1-2p_j}{1-p_j} 1_\Omega + \frac{1}{1-p_j} 1_{A_j}.
\end{equation*}
As $j \to \infty$, $p_j \to \frac{1}{3}$, so
\begin{align*}
    \lim_{j \to \infty} \int_\Omega f_j(x,y) \phi(x,y) \, dx \, dy &= \frac{1 - 2(1/3)}{1 - 1/3} \int_\Omega \phi(x,y) \, dx \, dy \\
    &+ \frac{1}{1 - 1/3} \left( \frac{1}{3} \int_\Omega \phi(x,y) \, dx \, dy \right) \\
    & = \int_\Omega \phi(x,y) \, dx \, dy.
\end{align*}
Thus, $F_j$ converges weakly to $G$.
\end{proof}

\end{document}
